% EGA V, English translation by P. Blass, 1987 (cote 24): archive cote 24, site page 12 (PDF page 13).
% typescript. Draft compared: none.
% AI-transcribed by Claude Opus 5.5 (claude-opus-5-5, adaptive thinking, effort high), Message Batches API, batch ega24b (msgbatch_01RGJJeHdR96U16ZWFuUjtDx), 2026-11-13; not yet checked by eye.
\scanpage{24/12}
\mspage{12}

\uline{Proposition 2.6}

Let us suppose that for every irreducible component $X_i$ of X
we have dim $\overline{f(X_i)} > 0$ i.e. $Y_{i\eta} \neq \emptyset$ or that f is an immersion and
dim $\overline{f(X)} > 0$ [slightly illegible)(Ask AG)].

Then we have dim $\underline{Y}_\eta$ = dim X - 1.

We are reduced to the case where X is irreducible due to 2.5.
By the very construction Y is defined starting from $X_{k(\eta)}$ as the
divisor of a section of an invertible module over $X_{k(\eta)}$ (being the
inverse image of $O_{\underline{P}}(1)$).  On the other hand $X_{k(\eta)}$ is irreducible
(because X is such and $k(\eta)$ is a pure transcendental extension of
k which one should have indicated at the beginning of the N°) and
$Y_\eta \neq X_{k(\eta)}$ since the image of $Y_\eta$ in X (contrary to that of $X_{k(\eta)}$,
which is faithfully flat over X) is not equal to X, indeed it does
not contain the closed points of X because of 2.3.  It follows
that dim $Y_\eta$ = dim $X_{k(\eta)}$ - 1 = dim X - 1 (reference needed for the
last equality.) \uline{q.e.d.}

\uline{Proposition 2.7}

Let F be a quasi coherent module over X, hence $G_\eta$ over $Y_\eta$.
Let $Z_i$ be the associated prime cycles of F such that dim $\overline{f(Z_i)} > 0$.
Let $Z_{i\eta}$ be the inverse image of $Z_i$ in $Y_\eta$ then the $Z_{i\eta}$ are exactly
all the prime cycles associated to $G_\eta$ also their relation of
inclusion are the same or among the $Z_i$.

The last assertion is contained in 2.4.  On the other hand
since $Y \to X$ is a projective fibration, thus flat with fibers $(S_1)$
and irreducible it follows from par. 3* that the associated prime

\margin{* Ask Grothendieck about this reference.}
